\documentclass[11pt]{article}
\usepackage{mathptmx}        % Times New Roman body + math
\usepackage[margin=1in]{geometry}
\usepackage{graphicx,amsmath,amssymb}

\setlength{\parskip}{4pt}
\title{Independent external validation of the SugarCode AI computational-biology platform against held-out public benchmarks}
\author{Udita Phookan}
\date{\today\\Independent research; computational work assisted by Instinct}
\begin{document}
\maketitle
\begin{abstract}
We subject SugarCode AI, a 95-module computational-biology platform, to an
independent external validation: every module family whose internal goldens
report verification against live data is re-tested here on held-out public
benchmarks its own tests never touched - fresh CRISPR guides mined from a
BRCA1 RefSeqGene segment disjoint from the module's training accessions and
scored against the CRISPOR reference implementation of the CFD off-target
model; a 500-variant stratified sample of the current ClinVar release; 4{,}956
single-point mutations from SKEMPI 2.0; the BiGG \emph{E. coli} core
metabolic model; and a globin-family alignment produced by EBI Clustal Omega.
We report per-module agreement metrics with bootstrap confidence intervals,
publish all fixtures with provenance, and document failure modes without
suppressing negative results.
\end{abstract}
\section{Introduction}
Computational-biology toolkits are typically validated against the same data
that informed their development, inflating confidence. Independent, held-out
validation against public reference implementations and current database
releases is the standard this platform must meet to be usable in laboratory
settings. This study asks a single question: does each verified claim in the
platform's own status ledger survive contact with data it has never seen?
We answer it module family by module family, with frozen fixtures, pinned
random seeds, and metrics reported with 95\% bootstrap confidence intervals.
All code, fixtures, and results are released in this repository.
\section{Related Work and Benchmarks}
\emph{Pending validation runs; filled on completion.}
\section{Methods}
For each module family we constructed a held-out test disjoint from every
fixture used during platform development. Negative results are preserved and
reported alongside positives; no test was re-designed after seeing its outcome.
\subsection{Data and Provenance}
All datasets are public and carry recorded provenance (see
\texttt{data/README.md}): NCBI ClinVar \texttt{variant\_summary.txt.gz}
(retrieved 2026-09-24; 500-variant stratified reservoir sample from
1{,}517{,}551 eligible GRCh38 SNVs with criteria-provided review status);
SKEMPI 2.0 \cite{skempi} (150-mutation frozen subset, seed 7); BiGG
\texttt{e\_coli\_core} \cite{bigg}; BRCA1 RefSeqGene NG\_005905.2:60000-70000
(NCBI efetch); CRISPOR reference distribution \cite{crispor}; 25 reviewed
globin-family UniProtKB sequences aligned with EBI Clustal Omega
\cite{clustalo}.
\subsection{Mathematical Formulation}
CFD off-target scoring follows Doench et al. 2016 \cite{doench2016}:
$\mathrm{CFD} = \prod_i m_i \cdot p(\mathrm{PAM})$, with per-position
mismatch penalties $m_i$ from the reference distribution. Binding free-energy
changes are $\Delta\Delta G = RT\ln(K_d^{mut}/K_d^{wt})$ at $T=300$ K
($RT = 0.596$ kcal/mol). Flux balance analysis solves
$\max_v c^\top v$ subject to $Sv = 0$, $l \le v \le u$, using the HiGHS
simplex solver; the published aerobic-glucose optimum for
\texttt{e\_coli\_core} is $0.8739\ \mathrm{h^{-1}}$ \cite{bigg}. Agreement is
quantified with Pearson $r$, Mann-Whitney AUC, and binary-classification
metrics; confidence intervals are percentile bootstrap with 1{,}000 resamples
at fixed seed.
\subsection{Implementation and Testing}
The validation suite is a pytest package (\texttt{sc\_validate}) with a
repo-wide no-stub audit: no fixture may contain placeholder data. Live-data
re-runs are separate scripts under \texttt{scripts/}; committed tests run
hermetically against frozen fixtures. CI runs the suite on every push.
\section{Results}
\emph{Pending validation runs; filled on completion.}
\section{Benchmark Comparison}
\emph{Pending validation runs; filled on completion.}
\section{Failure Modes and Negative Results}
\emph{Pending validation runs; filled on completion.}
\section{Discussion}
\emph{Pending validation runs; filled on completion.}
\section{Conclusion}
\emph{Pending validation runs; filled on completion.}
\bibliographystyle{plain}
\begin{thebibliography}{99}
\bibitem{doench2016} Doench JG et al. Optimized sgRNA design to maximize
activity and minimize off-target effects of CRISPR-Cas9. Nat Biotechnol
34:184-191, 2016.
\bibitem{crispor} Concordet JP, Haeussler M. CRISPOR: intuitive guide
selection for CRISPR/Cas9 genome editing experiments and screens. Nucleic
Acids Res 46:W242-W245, 2018.
\bibitem{skempi} Jankauskait\.e J et al. SKEMPI 2.0: an updated benchmark of
changes in protein-protein binding energy, kinetics and thermodynamics upon
mutation. Bioinformatics 35:462-469, 2019.
\bibitem{bigg} King ZA et al. BiGG Models: a platform for integrating,
standardizing and sharing genome-scale models. Nucleic Acids Res 44:D515-522,
2016.
\bibitem{clustalo} Sievers F et al. Fast, scalable generation of high-quality
protein multiple sequence alignments using Clustal Omega. Mol Syst Biol 7:539,
2011.
\bibitem{tavtigian} Tavtigian SV et al. Modeling the ACMG/AMP variant
classification guidelines as a Bayesian classification framework. Genet Med
20:1054-1060, 2018.
\end{thebibliography}
\end{document}
